\scanpage{119}
\mspage{118}


2) $\unsure{\otimes|\Omega}$-Functor $F:C\handwritten{\longrightarrow}C'$ (with $C,C'$ thingy-cate\rlap{t}gories).

The data are a \deleted{additive} functor $F$, plus \deleted{d}\unsure{\rlap{e}a|a}\rlap{a}n isomorphism\deleted{s} of \unsure{bifunctor|bifunctors}
\[
c(X,Y):F(X\otimes Y)\handwritten{\xrightarrow{\sim}}F(X)\otimes F(Y)
\]
\deleted{and isomorphism}
\[
\deleted{i:F(\underline{1}_C)\unsure{\text{xxx}|\longrightarrow}\underline{1}_{C'}}
\]
satisfying the conditions:

a) Compatibility \inserted{of $c$} with the datum of $\handwritten{\phi,\phi'}$

b) Compatibility \inserted{of $c$} with the datum of $\handwritten{\unsure{\mu|\psi},\mu'}$

c) \deleted{Compatibility of $i$ with and $c$}\inserted{Exists isomorphism\ \deleted{fonct} $i\handwritten{\unsure{\rlap{(}:|:}}F(\underline{1})\handwritten{\simeq}\underline{1}'$ compatible with $c$ and\handwritten{\begin{tikzpicture}[x=.65cm,y=.65cm,baseline=-.2ex]\draw (0,.28) .. controls (.3,.08) and (.77,.24) .. (.93,.05) .. controls (1.21,-.4) and (1.88,-.18) .. (2.26,.1)--(2.25,.7);\draw[line width=.15pt] (-.2,-.5) .. controls (-.35,-.13) and (-.34,.01) .. (-.12,.17);\end{tikzpicture}}} with the isomorphisms $u\inserted{u'\handwritten{\begin{tikzpicture}[x=.4cm,y=.4cm,baseline=-.2ex]\draw (0,-.25)--(0,.21) .. controls (.4,.48) and (.5,.83) .. (.75,.87) .. controls (.98,.9) and (1.08,1.09) .. (1.1,1.35);\end{tikzpicture}}}$ (or that \deleted{wh} which amounts to the same thing, with \inserted{$c$ and\handwritten{\begin{tikzpicture}[x=.45cm,y=.45cm,baseline=-.2ex]\draw (0,.7) .. controls (.25,.47) and (.72,.67) .. (.99,.35) .. controls (1.08,.16) and (1.14,-.08) .. (1.1,-.23);\draw (1.1,-.23)--(1.55,.88);\end{tikzpicture}}} $v,v'$).

NB This $i$ will then be uniquely determined.

d) $F$ is additive.

3) Homomorphism from a $\otimes$-functor $F$ to another $G$: it is a homomorphism of functors $f:F\handwritten{\longrightarrow}G$ satisfying th\rlap{e}e\deleted{s} condition

a) compatibility with \deleted{$c$, and} $c_F,c_G$

\handwritten{\deleted{\unsure{[a)|(A)} \unsure{for} \illegible{2}}
\begin{tikzpicture}[x=.65cm,y=.65cm,baseline=-.2ex]
\draw[line width=.2pt] (0,-.25) .. controls (.13,.1) and (.3,.08) .. (.48,.07)--(2.62,.21) .. controls (2.9,.25) and (3.08,.3) .. (3,.48);
\end{tikzpicture}}

\deleted{b)} \deleted{\unsure{(a)|(A)}}\handwritten{(}compatibility with $i_F,i_G$ (\deleted{it}\handwritten{\unsure{it|for}} is \deleted{perhaps}\inserted{\handwritten{\deleted{\unsure{probably}} probably not a}} formal consequence of a) - to look at\deleted{)}\handwritten{; \makebox[0pt][l]{\raisebox{1.4ex}{\hspace{.1em}\begin{tikzpicture}[x=1em,y=1em]\fill (0,0) circle (.035);\end{tikzpicture}}}cf below)}

\handwritten{4)} \uline{Rem\unsure{\rlap{\raisebox{1.2ex}{\hspace{.25em}\begin{tikzpicture}[x=1em,y=1em]\draw[line width=1pt] (0,0)--(0,.17);\end{tikzpicture}}}a|a}rk\handwritten{s}}\qquad To separate \inserted{\handwritten{what precedes\begin{tikzpicture}[x=.55cm,y=.55cm,baseline=-.2ex]\draw (0,-.35) .. controls (-.1,.21) and (.15,.41) .. (.48,.42)--(2.36,.51) .. controls (2.54,.53) and (2.55,.63) .. (2.5,.79);\end{tikzpicture}}} into several aspects of the datum of a bifunctor\handwritten{$\otimes$}:

a) Associativity datum, subject to the pentagon axiom. A functor ``compatible with \inserted{associat\deleted{t}ive\handwritten{\begin{tikzpicture}[x=.5cm,y=.5cm,baseline=-.2ex]\draw (0,-.3)--(0,.27) .. controls (.55,.36) and (1.56,.29) .. (2.15,.38) .. controls (2.4,.43) and (2.55,.62) .. (2.55,.82);\end{tikzpicture}}} $\otimes$'' from $C$ to $C'$ is a functor $+$ a $c$ , compatible with the datum of $\handwritten{\phi}$ .

b) Two-sided unit datum $\underline{1},u,v$ (unique up to unique isomorphism). A functor ``compatible with \inserted{\handwritten{$\otimes$ with\begin{tikzpicture}[x=.45cm,y=.45cm,baseline=-.2ex]\draw (0,-.25) .. controls (.28,.14) and (.51,.46) .. (.78,.51) .. controls (1.03,.57) and (1.32,.73) .. (1.4,.93);\end{tikzpicture}}} two-sided\deleted{s} unit\deleted{s}'' is an $F$ equipped with $c$, such that there exists $i$ compatible with $c,u,v$ and $c',u',v'$ (which $i$ will then be unique).

\margin{\handwritten{\uline{NB} The monoid of endomorphisms of $\underline{1}$ is then commutative.}}

\handwritten{c)} Commutativity datum $\handwritten{\mu}$ , satisfying the reflexivity condition.\deleted{(} A functor
